\chapter{Questionnaire Items}
\label{app:questionnaires}

Both instruments are reproduced item by item from the forms handed to participants, without their layout. Section~\ref{subsec:questionnaire-instruments} describes when each was administered and Section~\ref{subsec:statistical-framework} how each was scored.

\section{Pre-experiment questionnaire}

\paragraph{About you.} Age in years; gender (woman, man, prefer not to say); field of study or occupation; and how often the participant had interacted with a robot before the session (never, once or twice, a few times, many times, regularly).

\paragraph{Views on robots.} Fourteen statements about robots in general, each rated from 1 (strongly disagree) to 5 (strongly agree) \citep{Nomura2006NARS}.

\begin{enumerate}\setlength{\itemsep}{0pt}\setlength{\parskip}{0pt}
    \item I would feel uneasy if robots really had emotions.
    \item Something bad might happen if robots developed into living beings.
    \item I would feel relaxed talking with robots.
    \item I would feel uneasy if I was given a job where I had to use robots.
    \item If robots had emotions, I would be able to make friends with them.
    \item I feel comforted being with robots that have emotions.
    \item The word ``robot'' means nothing to me.
    \item I would feel nervous operating a robot in front of other people.
    \item I would hate the idea that robots or artificial intelligences were making judgements about things.
    \item I would feel very nervous just standing in front of a robot.
    \item I feel that if I depend on robots too much, something bad might happen.
    \item I would feel paranoid talking with a robot.
    \item I am concerned that robots would be a bad influence on children.
    \item I feel that in the future society will be dominated by robots.
\end{enumerate}

\section{Post-interaction questionnaire}

\paragraph{Impression of the robot.} Fifteen semantic differentials rated from 1 (left pole) to 5 (right pole) \citep{Bartneck2009Godspeed}. Items 1--5, 6--10 and 11--15 form the naturalness, animacy and likeability domains.

\begin{center}
\small
\begin{tabular}{@{}rll@{\hspace{2.5em}}rll@{}}
\toprule
1 & Fake & Natural & 9 & Inert & Interactive \\
2 & Machinelike & Humanlike & 10 & Apathetic & Responsive \\
3 & Unconscious & Conscious & 11 & Dislike & Like \\
4 & Artificial & Lifelike & 12 & Unfriendly & Friendly \\
5 & Moving rigidly & Moving elegantly & 13 & Unkind & Kind \\
6 & Dead & Alive & 14 & Unpleasant & Pleasant \\
7 & Stagnant & Lively & 15 & Awful & Nice \\
8 & Mechanical & Organic & & & \\
\bottomrule
\end{tabular}
\end{center}

\paragraph{Words associated with the robot.} Twelve words rated from 1 (not associated at all) through 5 (moderately associated) to 9 (strongly associated) \citep{Carpinella2017RoSAS}: happy, scary, feeling, strange, social, awful, organic, awkward, compassionate, dangerous, emotional, aggressive.

\paragraph{Experience of the robot.} Four statements rated from 1 (strongly disagree) to 5 (strongly agree). (1) The robot's hunger expression felt natural. (2) The robot's hunger felt like a genuine internal state rather than a performance. (3) The robot felt like a part of the environment. (4) The robot's hunger made me want to help it.

\paragraph{In your own words.} Four open prompts: (1) How would you describe your overall experience of interacting with the robot? (2) Did the robot's behaviour feel natural or unnatural to you? What gave you that impression? (3) Was there a particular moment that stood out to you, whether positive or negative? (4) Is there anything else you would like to tell us about the robot or this session?
